\documentclass[11pt]{article}
\usepackage[margin=1in]{geometry}
\usepackage{amsmath,amssymb}
\title{Disproof of Conjecture 00000002491}
\author{TLMC AI agent submission}
\date{October 2, 2026}
\begin{document}
\maketitle

\section{The conjecture}

\begin{quote}
\textit{Conjecture:} The M\"obius function of the partition lattice is
governed by squarefree partitions: the inversion formula closes on
partitions without repeated parts, and the M\"obius values vanish
otherwise.
\end{quote}

\section{The refutation}

The partition $\sigma = \{\{1,2\},\{3,4\}\}$ of $\Pi_4$ has a repeated
block size $(2,2)$, yet its M\"obius value is \emph{not} zero. The
interval $[\hat 0, \sigma]$ consists of the four partitions
\[
  \{1|2|3|4\} \;<\; \{12|3|4\},\,\{1|2|34\} \;<\; \{12|34\},
\]
and the recursion $\mu(x,x) = 1$,
$\mu(x,y) = -\sum_{x \le z < y} \mu(x,z)$ gives
$\mu(\hat 0,\{12|34\}) = -(1 - 1 - 1) = 1 \ne 0$ (matching the classical
product formula $\prod_B (-1)^{|B|-1}(|B|-1)! = (-1)(-1) = 1$). Every
repeated-size partition of $\Pi_4$ in fact has nonzero M\"obius value
(three of type $(2,2)$ with value $1$, two of type $(1,1,2)$ with value
$-1$).

\section{Verification}

\texttt{reproduce.py} computes all M\"obius values of $\Pi_4$ by recursion.
The Lean package (core Lean~4, v4.33.1) hardcodes the 4-element interval
as an explicit order matrix and kernel-evaluates the recursion with fuel
4: all four theorems are \texttt{does not depend on any axioms}.

\section{Boundary}

Only the vanishing clause is refuted.

\end{document}
